\section{Original marginals in regional coordinates}

Let $V$ be a finite set of sites with finite local coordinate sets $\mathcal B_v$.
For $T\subseteq V$, write $\mathcal B_T=\prod_{v\in T}\mathcal B_v$ and
$\mathcal H_T=\mathbb C^{\mathcal B_T}$. Let $\Omega$ be an arbitrary vector
in $\mathcal H_V$, and let
$\rho_T=\operatorname{Tr}_{T^c}|\Omega\rangle\langle\Omega|$.
The coordinate sets may be empty, and no normalization is imposed on $\Omega$.

\begin{definition}[The three original physical regions]
  \label{def:regional_physical_regions}
  \lean{FiniteProduct.regionalPhysicalRegion}
  For regions $Q,Y\subseteq V$, define
  \begin{align*}
    T_0=Q,\qquad T_1=Y,\qquad T_2=(Q\cup Y)^c.
  \end{align*}
\end{definition}

\begin{theorem}[The original reductions after regional grouping]
  \label{thm:regional_physical_singleton_densities}
  \lean{Matrix.regional_physicalSingletonDensity_eq_reducedPure}
  \uses{def:regional_physical_regions,def:regional_five_factor_coordinates,def:physical_singleton_density}
  Suppose $Q$ and $Y$ are disjoint, and fix any two finite auxiliary
  coordinate sets $C,R$. Let $e_{Q,Y}$ be restriction of configurations to
  $Q$, $Y$, and $(Q\cup Y)^c$, in that order, and put
  $\psi=\Omega\circ e_{Q,Y}^{-1}$. In the five-factor coordinates
  $(\mathcal B_Q,\mathcal B_Y,\mathcal B_{(Q\cup Y)^c},C,R)$, the three
  physical singleton reductions of $|\psi\rangle\langle\psi|$ satisfy
  \begin{align*}
    \rho^{\psi}_0=\rho_Q,\qquad
    \rho^{\psi}_1=\rho_Y,\qquad
    \rho^{\psi}_2=\rho_{(Q\cup Y)^c}.
  \end{align*}
  These are the original physical marginals used in the component estimates
  of \cite[Section~8]{OpenAI2026AreaLaw}.
\end{theorem}

\begin{proof}
  For the first reduction, the discarded coordinates are
  $\mathcal B_Y\times\mathcal B_{(Q\cup Y)^c}$, which are the original
  coordinates on $Q^c$. For the second, they are
  $\mathcal B_Q\times\mathcal B_{(Q\cup Y)^c}$, the coordinates on $Y^c$.
  For the third, they are $\mathcal B_Q\times\mathcal B_Y$, the coordinates
  on $Q\cup Y$. In each case the retained coordinate is unchanged.
  Reindexing the finite sum defining the partial trace proves the equality.
\end{proof}

\begin{theorem}[Preservation of entropy and surprisal moments]
  \label{thm:regional_physical_spectral_statistics}
  \lean{Matrix.regional_physicalSingleton_spectral_statistics}
  \uses{thm:regional_physical_singleton_densities}
  Under the hypotheses of Theorem~\ref{thm:regional_physical_singleton_densities},
  let $j\in\{0,1,2\}$. For a positive semidefinite matrix $\rho$ with
  eigenvalues $\lambda_a$, define
  \begin{align*}
    S(\rho)&=-\sum_a\lambda_a\log\lambda_a,\qquad
    M_\rho(v)=\sum_a\lambda_a\exp(-v\log\lambda_a).
  \end{align*}
  A zero eigenvalue contributes zero to both expressions. For every real $v$,
  \begin{align*}
    S(\rho^{\psi}_j)&=S(\rho_{T_j}),\qquad
    M_{\rho^{\psi}_j}(v)=M_{\rho_{T_j}}(v)
  \end{align*}
\end{theorem}

\begin{proof}
  The two matrices are equal in the same original regional basis by
  Theorem~\ref{thm:regional_physical_singleton_densities}. Their eigenvalue
  lists therefore agree, and so do the displayed finite sums.
\end{proof}
